\part{Cifrari simmetrici}

\newpage

\section{Advanced Encryption Standard}

Nel 1997 il NIST indisse un concorso pubblico per selezionare un nuovo standard di cifratura simmetrica destinato a
sostituire il DES, ormai ritenuto non più adeguato dal punto di vista della sicurezza. Il DES, sviluppato negli anni
'70, utilizzava infatti una chiave di soli 56 bit, rivelatasi vulnerabile ad attacchi di brute force con l'aumento
della potenza di calcolo disponibile.

Il concorso si concluse nel 2000 con la selezione dell'algoritmo Rijndael, progettato dai crittografi belgi Joan Daemen
e Vincent Rijmen, che venne successivamente adottato come Advanced Encryption Standard (AES).

Le proposte candidate furono valutate principalmente sulla base di tre criteri: la sicurezza, considerata essenziale;
l'efficienza computazionale, intesa in termini di velocità e requisiti di memoria in diversi contesti implementativi; e
le caratteristiche algoritmiche e implementative, quali semplicità, flessibilità e facilità di adozione.

AES fu scelto per la sua combinazione di sicurezza, efficienza, prestazioni e flessibilità. La sua struttura a blocchi,
il supporto a chiavi di lunghezza diversa e la resistenza ai principali attacchi crittoanalitici lo resero il candidato
ideale per il nuovo standard.

\subsection{Struttura di AES}

\subsubsection{Struttura generale}

AES è un cifrario simmetrico che opera su blocchi di lunghezza fissa pari a 128 bit, cioè 16 byte. L'algoritmo ammette
tre possibili lunghezze di chiave: 128 bit, 192 bit e 256 bit. In corrispondenza di tali lunghezze si parla
rispettivamente di AES-128, AES-192 e AES-256. Dal punto di vista strutturale, AES è un \emph{cifrario iterato}: il
blocco di input viene sottoposto a una sequenza di trasformazioni elementari che si ripetono round dopo round.
Indicando con $N$ il numero di round, si ha $N = 10$, $N = 12$ oppure $N = 14$ a seconda che la chiave abbia lunghezza
rispettivamente di 128, 192 o 256 bit.

Tutte le operazioni agiscono su una rappresentazione interna del blocco, chiamata \emph{State}, costituita da una
matrice quadrata $4 \times $4 di byte. In altre parole, i 16 byte del plaintext non vengono trattati come una semplice
sequenza lineare, ma come una struttura matriciale che viene modificata progressivamente nel corso della cifratura.
L'organizzazione dei dati in ingresso, della matrice di stato e dell'output è illustrata in
\figsubref{fig:aes-data-structure}{a}

L'ordinamento dei byte all'interno dello \emph{State} avviene per colonne. Ciò significa che i primi quattro byte del
blocco occupano la prima colonna della matrice, i quattro byte successivi la seconda colonna, e così via. Questa
convenzione è importante, perché tutte le trasformazioni interne di AES, in particolare quelle che agiscono sulle righe
o sulle colonne della matrice, sono definite rispetto a tale disposizione.

Anche la chiave iniziale viene organizzata in forma matriciale e successivamente sottoposta a un processo di espansione
(\emph{key expansion} o \emph{key schedule}), mediante il quale si ottengono le \emph{round keys} utilizzate nei
diversi round. In questo senso, la chiave non interviene sempre in modo identico durante l'esecuzione dell'algoritmo,
ma viene preliminarmente trasformata in una sequenza di sottochiavi, una per ciascun round, più una sottochiave
iniziale, come mostrato in \figsubref{fig:aes-data-structure}{b}.

Ad alto livello, il funzionamento di AES può essere descritto come mostrato in
\figsubref{fig:aes-encryption-process}{}. Dato un plaintext $x$, esso viene inizialmente copiato nello \emph{State}.
Prima dell'inizio dei round veri e propri, si applica una trasformazione iniziale di \texttt{AddRoundKey}, che combina
lo \emph{State} con la prima round key; questa trasformazione iniziale può essere considerata a tutti gli effetti come
un round $0$. Seguono poi i successivi $N - 1$ round, ciascuno dei quali è composto da quattro trasformazioni:
\texttt{SubBytes}, \texttt{ShiftRows}, \texttt{MixColumns} e \texttt{AddRoundKey}. L'ultimo round ha invece una
struttura leggermente diversa, poiché non include la trasformazione \texttt{MixColumns}; esso è composto soltanto da
\texttt{SubBytes}, \texttt{ShiftRows} e \texttt{AddRoundKey}. Tutte queste operazioni ricevono in input una matrice $4
\times 4$ di byte e producono in output una nuova matrice della stessa forma. Dopo l'ultimo round, la matrice di stato
viene copiata in una matrice di output, che costituisce il ciphertext.

Questa organizzazione mostra che AES possiede una struttura regolare e fortemente modulare. Ogni round combina infatti
una fase di sostituzione non lineare, una fase di permutazione delle posizioni, una trasformazione lineare sulle
colonne e un'operazione di mescolamento con la sottochiave di round.

\begin{figure}[H]
  \centering
  \includegraphics[width=0.66\linewidth]{2_aes_structure.jpg}
  \vspace{-4mm}
  \caption{AES Encryption Process} \label{fig:aes-encryption-process}
\end{figure}

\vspace{-4mm}

\begin{figure}[H]
  \centering
  \includegraphics[width=0.975\linewidth]{2_aes_data_structure.jpg}
  \vspace{-3mm}
  \caption{AES Data Structure} \label{fig:aes-data-structure}
\end{figure}

\subsubsection{Struttura dettagliata}

La \figsubref{fig:aes-encryption-decryption}{a} mostra il cifrario AES più in dettaglio, indicando la sequenza di
trasformazioni in ciascun round e mostrando la corrispondente operazione di decryption.

\begin{figure}[H]
  \centering
  \includegraphics[width=0.7\linewidth]{2_aes_encryption_decryption.jpg}
  \vspace{-2mm}
  \caption{AES Encryption and Decryption} \label{fig:aes-encryption-decryption}
\end{figure}

Una prima osservazione importante è che AES non appartiene alla famiglia dei cifrari di Feistel. In una struttura di
Feistel, infatti, il blocco viene suddiviso in due metà, una delle quali viene usata per modificare l'altra, seguita
poi da uno scambio delle due parti. In AES, invece, l'intero blocco viene trattato in ogni round come un'unica matrice
di byte, sulla quale agiscono trasformazioni di sostituzione, permutazione e mescolamento lineare. Per questo motivo,
AES può essere ricondotto al paradigma delle \emph{Substitution-Permutation Networks} (SPN), pur presentando una
struttura specifica propria.

\begin{figure}[H]
  \centering
  \includegraphics[width=0.85\linewidth]{2_aes_encryption_round.jpg}
  \vspace{-2mm}
  \caption{AES Encryption Round} \label{fig:aes-encryption-round}
\end{figure}

\subsection{Funzioni di trasformazione}

\subsubsection{Substitute Bytes}

\subsubsection{Shift Rows}

\subsubsection{Mix Columns}

\subsubsection{Add Round Key}

\subsection{Key Expansion}

\newpage
\part{Matematica}

